% ============================================================
% Proposition 4.5: Compensation Analysis
% Status: FULL RECONSTRUCTION
% ============================================================

% --- Definitions and Assumptions ---
\begin{definition}[Intersection Discovery Rate]
The rate at which agent $A$ discovers compositional intersections
between domains $d_1$ and $d_2$ is:
\[
\Psi_A(d_1, d_2, t) = \Psi_0 \cdot \phi(d_1, d_2) \cdot
\sqrt{q_A(d_1, t) \cdot q_A(d_2, t)},
\]
as defined in Equation~\eqref{eq:discovery-rate}, where
$q_A(d, t) = \sigma_A(d, t) \cdot \delta_A^{\text{relative}}(d, t)$
is the mastery quality \eqref{eq:mastery-quality}.
\end{definition}

\begin{definition}[Attentional Fidelity ODE]
During training, $\alpha_A$ evolves according to
Equation~\eqref{eq:alpha-ode}:
\[
\dot{\alpha}_A = \gamma C_A (1 - \alpha_A) - \zeta_\alpha \alpha_A R_A^{\text{surface}}.
\]
The coupling to $\sigma_A$ occurs through the $\sigma_A$ ODE
\eqref{eq:sigma-ode}:
\[
\dot{\sigma}_A = \sigma_A \cdot
[\rho P_A \alpha_A (1 - \gamma_\sigma \sigma_A) - \epsilon_\sigma \Omega_{SL}],
\]
where $\alpha_A$ appears as a multiplicative gate on the growth term.
\end{definition}

\begin{assumption}[Evaluation-Time Stationarity]
At evaluation time, $\sigma_A$ and $\alpha_A$ are considered fixed
(measured at the end of training). The dynamics that produced them
are not relevant to the evaluation-time expression for $\Psi_A$.
\end{assumption}

\begin{assumption}[No $\alpha_A$ Dependence in $\Psi_A$]
The evaluation-time expression \eqref{eq:discovery-rate} contains
no $\alpha_A$ term. This follows from the definition of $q_A$ as
$\sigma_A \cdot \delta_A^{\text{relative}}$, which depends only on
$\sigma_A$ and $\delta_A$.
\end{assumption}

% --- Proposition Statement ---
\begin{proposition}[Compensation Analysis]
\label{prop:compensation}
High attentional fidelity $\alpha_A$ cannot compensate for low schema
coherence $\sigma_A$ in the intersection discovery rate:
\[
\Psi_A \propto \sigma_A \cdot \delta_A^{\text{relative}}
\quad \text{(independent of } \alpha_A \text{ at evaluation time)}.
\]
While $\alpha_A$ gates the \emph{growth} of $\sigma_A$ during training
\eqref{eq:sigma-ode}, it does not appear in the evaluation-time
expression for $\Psi_A$. Thus, a benchmark measuring intersection
discovery is insensitive to $\alpha_A$ once $\sigma_A$ is fixed.
\end{proposition}

% --- Proof ---
\begin{proof}
We analyze the evaluation-time expression and the training-time
dynamics separately.

\paragraph{Step 1: Evaluation-time independence.}
By Equation~\eqref{eq:discovery-rate} and
\eqref{eq:mastery-quality}:
\[
\Psi_A(d_1, d_2, t) = \Psi_0 \cdot \phi(d_1, d_2) \cdot
\sqrt{(\sigma_A(d_1) \cdot \delta_A^{\text{relative}}(d_1))
\cdot (\sigma_A(d_2) \cdot \delta_A^{\text{relative}}(d_2))}.
\]
Taking the partial derivative with respect to $\alpha_A$:
\[
\frac{\partial \Psi_A}{\partial \alpha_A} = 0,
\]
because $\alpha_A$ does not appear anywhere in the expression.
The same applies to the joint sensitivity:
\[
\frac{\partial^2 \Psi_A}{\partial \sigma_A \partial \alpha_A} = 0.
\]
Thus, for fixed $\sigma_A$ and $\delta_A^{\text{relative}}$, $\Psi_A$
is independent of $\alpha_A$. No amount of attentional fidelity can
directly increase the intersection discovery rate if $\sigma_A$ is low.

\paragraph{Step 2: Training-time dependence.}
The ODE \eqref{eq:sigma-ode} shows that $\alpha_A$ gates the growth
rate of $\sigma_A$:
\[
\dot{\sigma}_A = \sigma_A \cdot
[\rho P_A \alpha_A (1 - \gamma_\sigma \sigma_A) - \epsilon_\sigma \Omega_{SL}].
\]
The term $\rho P_A \alpha_A$ is the growth engine. When $\alpha_A = 0$:
\[
\dot{\sigma}_A = \sigma_A \cdot [0 - \epsilon_\sigma \Omega_{SL}]
= -\sigma_A \epsilon_\sigma \Omega_{SL} \leq 0.
\]
Schema coherence decreases monotonically when attention is absent.
When $\alpha_A > 0$, the growth term provides upward pressure on
$\sigma_A$ (when the bracket is positive).

The result is that $\alpha_A$ determines the \emph{trajectory} of
$\sigma_A$ during training, but not the evaluation-time expression
of $\Psi_A$. Once training is complete and $\sigma_A$ is fixed,
$\Psi_A$ is determined by $\sigma_A$ alone (together with
$\delta_A^{\text{relative}}$), independently of $\alpha_A$.

\paragraph{Step 3: Compensatory relationships that fail.}
Consider an agent with low $\sigma_A$ but high $\alpha_A$. Can the
high $\alpha_A$ compensate for low $\sigma_A$ in $\Psi_A$? The
evaluation-time expression:
\[
\Psi_A \propto \sqrt{\sigma_A \cdot \delta_A^{\text{relative}}}.
\]
$\alpha_A$ does not appear. Compensation is impossible in the
intersection discovery rate.

Consider an agent with high $\alpha_A$ during training. This would
(given sufficient $\rho$, $P_A$, and low $\Omega_{SL}$) produce higher
$\sigma_A$ at the end of training. But this is an indirect effect
through the dynamics, not a direct compensation mechanism. If other
factors ($\rho P_A$ small, $\Omega_{SL}$ large) prevent $\sigma_A$
from growing despite high $\alpha_A$, the high $\alpha_A$ is
ineffective.

\paragraph{Step 4: Implication for benchmark design.}
A benchmark measuring intersection discovery (e.g., compositional
generalisation tests) assesses the product of $\sigma_A$ and
$\delta_A^{\text{relative}}$. It does not assess $\alpha_A$ directly.
To measure $\alpha_A$, a benchmark must probe attention during
training (e.g., through distractor tasks or attention-weight
analysis). The $\Sigma$-Model predicts that discovery-rate benchmarks
and attention benchmarks are complementary: they measure different
latent variables that interact during training but factor at
evaluation time. $\square$
\end{proof}

% --- Boundary Cases ---
\begin{boundary-case}
\textbf{$\sigma_A = 0$ (no schema coherence):}
$\Psi_A = 0$ regardless of $\alpha_A$, $\delta_A^{\text{relative}}$,
or any other parameter. The intersection discovery rate is identically
zero. No compensation is possible.

\textbf{$\alpha_A = 0$ but $\sigma_A > 0$ (attention collapse with existing coherence):}
If $\sigma_A$ was built before attention collapsed, the evaluation-time
$\Psi_A$ remains proportional to $\sqrt{\sigma_A \delta_A}$ even with
$\alpha_A = 0$. However, $\sigma_A$ will decay under the dynamics
($\dot{\sigma}_A \leq 0$ when $\alpha_A = 0$), so this is a transient
state.

\textbf{$\delta_A^{\text{relative}} = 0$ (zero depth):}
$\Psi_A = 0$ regardless of $\sigma_A$ or $\alpha_A$. Depth is
necessary for intersection discovery — even perfect schema coherence
cannot compensate for lack of knowledge. This is the ``no free lunch''
boundary.

\textbf{Both $\sigma_A$ and $\alpha_A$ are high:}
The full $\Psi_A$ expression is evaluated with $\sigma_A \approx 1$,
giving $\Psi_A \propto \sqrt{\delta_A^{\text{relative}}}$. The
intersection discovery rate is limited only by depth.
\end{boundary-case}

% --- Circular Dependency Check ---
\begin{circularity-check}
\textbf{Dependencies on earlier results:} The proof uses the
definitions of $\Psi_A$ \eqref{eq:discovery-rate},
$q_A$ \eqref{eq:mastery-quality}, and the $\sigma_A$ ODE
\eqref{eq:sigma-ode}. It does not depend on any specific proposition.

\textbf{Forward dependencies:} None. This is a self-contained
analysis of the evaluation-time expression.

\textbf{Self-circularity:} The claim that ``$\alpha_A$ does not
appear in $\Psi_A$'' is definitionally true given the model equations.
The proof demonstrates the absence of $\alpha_A$ from $\Psi_A$ by
inspection of the definitions. This is not a deep mathematical
result but an architectural observation about the $\Sigma$-Model.
The significance lies in the implications for benchmark design and
interpretability, not in the formal difficulty of the proof. The
analysis of the training-time dynamics ($\alpha_A$ gates $\sigma_A$
growth) is also definitional given the ODE structure. The proposition
is best understood as a formal statement of a non-trivial architectural
constraint rather than a non-obvious mathematical fact.
\end{circularity-check}
